% N14 methods-only insertion.  Requires the amsmath/amssymb/bm packages
% already loaded by main_stage_draft.tex.  New citation keys are listed in
% INTEGRATION_NOTES.md.  No selection or confirmation outcomes appear here.
\section{Causal re-tensorization of a fixed RFF controller}
\label{sec:n14_method}

\subsection{Coefficient-guided coordinate layout}
The fixed Gaussian map in Eq.~(4) uses $\ell_r=3.9$ and supplies $D=500$
features throughout.  We use
$D_1=25$, $D_2=20$ and the sequential filtered-x MCC factor updates in
Eqs.~(11)--(12), with $\mu_a=\mu_b=\mu/2$, $\sigma_c=2$, and
$\varepsilon=10^{-8}$.  The initial active controller has four Kronecker
terms.  A structural proposal changes the assignment of these same 500
features to the $20\times25$ coefficient matrix; it does not select or
discard RFFs.

Let $\boldsymbol\pi$ list the original feature indices in the current
controller order, and let $\mathbf{P}_{\boldsymbol\pi}$ be its permutation
matrix.  Within a branch, both the output features and filtered-x features
use that order:
\begin{equation}
 \mathbf{z}^{\boldsymbol\pi}_n=\mathbf{P}_{\boldsymbol\pi}\mathbf{z}_n,
 \qquad
 \mathbf{q}^{\boldsymbol\pi}_n=\mathbf{P}_{\boldsymbol\pi}\mathbf{q}_n,
 \qquad
 \mathbf{w}^{\boldsymbol\pi}_n
   =\mathrm{vec}(\mathbf{B}_n\mathbf{A}_n^T).
 \label{eq:n14_permuted_map}
\end{equation}
At a proposal sample $t$, after the active branch's update on that sample,
we stably sort the \emph{signed} entries of its current coefficient vector:
\begin{equation}
 \mathbf{p}_t=\operatorname{argsort}_{\rm stable}
       (\mathbf{w}^{\boldsymbol\pi}_t),\qquad
 \boldsymbol\pi^+=\boldsymbol\pi[\mathbf{p}_t],\qquad
 \widetilde{\mathbf W}_t
   =\operatorname{mat}_{D_2\times D_1}
       (\mathbf{P}_{\mathbf p_t}\mathbf{w}^{\boldsymbol\pi}_t).
 \label{eq:n14_sort}
\end{equation}
Here $\operatorname{mat}_{D_2\times D_1}$ reverses columnwise
vectorization.  Simultaneously permuting the complete coefficients and the
features preserves every untruncated output exactly:
\begin{equation}
 (\mathbf{w}^{\boldsymbol\pi}_t)^T
       \mathbf{z}^{\boldsymbol\pi}_n
 = (\mathbf{P}_{\mathbf p_t}\mathbf{w}^{\boldsymbol\pi}_t)^T
       (\mathbf{P}_{\mathbf p_t}\mathbf{z}^{\boldsymbol\pi}_n).
 \label{eq:n14_invariance}
\end{equation}
The permutation is fixed for the lifetime of the proposed branch, and the
same permutation is applied to $\mathbf q_n$.  This identity describes the
full vector \emph{before} rank truncation; it is not a claim that a rank-two
candidate has identical output.

We compute a counted QR--Jacobi SVD
$\widetilde{\mathbf W}_t=\mathbf U\boldsymbol\Sigma\mathbf V^T$ and form
the best rank-two Frobenius approximation:
\begin{equation}
 \mathbf{B}_{c,t}=\mathbf U_{:,1:2}
          \boldsymbol\Sigma_{1:2}^{1/2},\qquad
 \mathbf{A}_{c,t}=\mathbf V_{:,1:2}
          \boldsymbol\Sigma_{1:2}^{1/2},\qquad
 \eta_{2,t}=
 \left(\frac{\sum_{j>2}\sigma_j^2}{\sum_j\sigma_j^2}\right)^{1/2}.
 \label{eq:n14_truncation}
\end{equation}
The candidate inherits the active branch's secondary-output FIR history.
If the decomposition fails to converge or $\eta_{2,t}>0.15$, the proposal
is rejected after charging its incurred computation.  Otherwise the
candidate continues the same sequential MCC updates as the active branch,
using its own permutation and factor states.

The first proposal follows sample $t=20{,}000$ (zero-based indexing) and
is the sole possible reduction from $R=4$ to $R=2$.  After each accepted
proposal, the next opportunity occurs 50{,}000 newly observed samples
later.  These later proposals start from the \emph{current} $R=2$
coefficients and produce an $R=2$ candidate in a newly sorted layout:
they update the coordinate arrangement and factors, not the number of
Kronecker terms.  A rejected or spectrally vetoed proposal ends structural
search for that run.  Neither a known change point nor a future coefficient
enters the proposal rule.

\subsection{Virtual secondary-path validation and actuator transition}
At each sample there is one physical actuator drive $y_n^{\rm act}$ and
one measured microphone error $e_n$.  With the known secondary-path model
$\widehat{\mathbf s}=\mathbf s$ in the present simulation, the observable
primary-plus-measurement disturbance is reconstructed as
\begin{equation}
 \widetilde d_n=e_n+
       \sum_{\ell=0}^{L_s-1}\widehat s_\ell
                              y_{n-\ell}^{\rm act}.
 \label{eq:n14_d_reconstruction}
\end{equation}
The active and candidate branches maintain separate hypothetical actuator
histories.  Filtering them with $\widehat{\mathbf s}$ gives
$\widehat y_{s,a,n}$ and $\widehat y_{s,c,n}$, respectively; their virtual
secondary-path residuals are
\begin{equation}
 \widetilde e_{a,n}=\widetilde d_n-\widehat y_{s,a,n},
 \qquad
 \widetilde e_{c,n}=\widetilde d_n-\widehat y_{s,c,n}.
 \label{eq:n14_virtual_residuals}
\end{equation}
The candidate is evaluated from the actual measured error and the
secondary-path model, but does not itself drive the actuator until
acceptance.  Each branch updates on $\widetilde d_n$ through
Eqs.~(11)--(12).  The MCC influence $\psi$ belongs to those factor
updates; it is not the candidate-selection loss.

The proposal is followed by 2{,}000 strictly subsequent samples of
candidate training and then 2{,}000 new samples of prequential validation.
For the validation set $\mathcal V_t$, let $s_n$ be a causal slow scale of
the actual absolute microphone error and set
$c_n=\max\{3s_n,10^{-3}\}$.  The gate uses clipped absolute virtual
residuals,
\begin{equation}
 L_j(t)=\sum_{n\in\mathcal V_t}
       \min\{|\widetilde e_{j,n}|,c_n\},
 \quad j\in\{a,c\},\qquad
 \frac{L_c(t)}{\max\{L_a(t),10^{-12}\}}
       \leq 10^{0.5/20}.
 \label{eq:n14_gate}
\end{equation}
The right-hand side permits at most a 0.5-dB increase in this
\emph{validation amplitude} score.  The cap is formed only from current
and past physical microphone errors.  Specifically, starting at zero,
the implementation uses
$r_n=|e_n|$ when $s_{n-1}=0$, and otherwise
$r_n=\min\{|e_n|,3\max(s_{n-1},10^{-3})\}$;
$f_n=0.98f_{n-1}+0.02r_n$.
A hold counter increments when
$n>1000$ and $f_n>1.28s_{n-1}$, and resets otherwise.
For hold count below 250,
$s_n=0.9995s_{n-1}+0.0005r_n$; otherwise
$s_n=s_{n-1}$.  This scale supplies the cap in Eq.~\eqref{eq:n14_gate}
and does not trigger rank growth in the present method.

Upon acceptance, the new branch becomes active, while the retiring
branch supplies a 100-sample actuator transition.  For transition sample
$j=1,\ldots,100$,
\begin{equation}
 y_{t+j}^{\rm act}
   =\left(1-\frac{j}{100}\right)y_{o,t+j}
       +\frac{j}{100}y_{c,t+j}.
 \label{eq:n14_ramp}
\end{equation}
The physical secondary path filters this \emph{single mixed output} with
its continuous FIR history.  The retiring and new branches retain their
own hypothetical histories for the transition; the new active branch is
then synchronized to the actual actuator history.  Candidate training,
validation, vetoes, rejection, and a candidate still pending at the record
endpoint all contribute to the arithmetic bill.  The a posteriori
factor-step contraction of Proposition~1 applies to the fixed-layout
updates in Eqs.~(11)--(12); it makes no contraction assertion about SVD
truncation or the actuator transition.

\subsection{Adaptive state and complete search arithmetic}
The active factor state contains $R(D_1+D_2)=45R$ scalar coefficients:
180 before the first accepted reduction and 90 thereafter.  All 500 RFFs
and their filtered histories remain in use.  A feature permutation stores
500 indices (approximately 1, 2, or 4 kB with 16-, 32-, or 64-bit indices);
overlapping active, candidate, and retiring branches may each need their
own permutation and factor state.

For $M=20$, $D=500$, and the conservative four-tap filtered-x count, the
fixed-rank cost including the physical secondary FIR is
\begin{equation}
 \begin{aligned}
 C_R&=12{,}512+1{,}655R
       &&\text{multiplications/sample},\\
 C_2&=15{,}822,\qquad C_4=19{,}132.
 \end{aligned}
 \label{eq:n14_fixed_cost}
\end{equation}
The four multiplications for the actual physical FIR are separate from
the fixed-controller count in the table immediately preceding this
section.  In particular, its 19{,}128 for $R=4$ becomes
19{,}132 under Eq.~\eqref{eq:n14_fixed_cost}.  At sample $n$ the dynamic
bill is computed from non-overlapping terms,
\begin{equation}
 \begin{aligned}
 C_n={}&(12{,}508+1{,}655R_{a,n})
       +I_{c,n}(12+1{,}655R_{c,n})\\
      &+I_{o,n}(12+1{,}655R_{o,n})
       +C_{\mathrm{manage},n}+C_{\mathrm{reorder},n}+4.
 \end{aligned}
 \label{eq:n14_dynamic_cost}
\end{equation}
The shared RFF projection, scaling, and filtered-x computation are charged
once even while branches overlap.  $I_c$ and $I_o$ indicate a live
candidate and retiring branch.  Management counts the causal scale and
phase logic (12 multiplications per nonfixed sample), private secondary
histories, four validation-score multiplications per validation sample,
two crossfade multiplications per ramp sample, and proposal factor-scale
work.  The one-time reorder bill includes coefficient assembly, the
counted QR and Jacobi sweeps, rank truncation, and the spectral-tail
calculation.  Rejected and unfinished proposals retain their incurred
terms in $\sum_{n=0}^{T-1}C_n/T$.  The frozen comparison is to an
equal-input, equal-initialization continuously trained $R=4$ controller;
the prespecified scalar-multiplication target is
$T^{-1}\sum_n C_n/C_4\leq 0.90$.  The offline calibration grid comprises
80 dynamic and 720 fixed-controller case--method--run evaluations; those
evaluations are reported separately from this per-run online bill.
Sorting comparisons, permutation memory
traffic, cosine and exponential evaluation, division, and wall-clock time
are separate resource measures; a reduction in Eq.~\eqref{eq:n14_dynamic_cost}
alone is not a hardware speedup claim.

\subsection{Relation to prior methods}
RFF correntropy ANC and Kronecker filtered-x MCC adaptation already have
close precedents \cite{ref11,n14_nkp_fxnsaf_mcc}; model-based virtual
candidate evaluation and structural selection have also been used in
nonlinear ANC \cite{n14_spiriti}.  Sorting or reordering weights to expose
low-rank tensor structure is likewise established outside ANC
\cite{n14_einsort,n14_neukron}.  The method studied here combines a
\emph{causal, coefficient-derived permutation of a fixed dense RFF map}
with synchronized feature and filtered-x coordinates, a single validated
$R=4\!\to\!2$ reduction, and subsequent validated $R=2\!\to\!2$
re-tensorization in an operating filtered-x ANC loop.  Its candidate
gate is the clipped virtual secondary-path score of
Eq.~\eqref{eq:n14_gate}, while its learning recursion remains MCC and
its online search operations enter Eq.~\eqref{eq:n14_dynamic_cost}.
